% AER-Article.tex for AEA last revised 22 June 2011
\documentclass[Department of Economics]{AEA}

% The mathtime package uses a Times font instead of Computer Modern.
% Uncomment the line below if you wish to use the mathtime package:
%\usepackage[cmbold]{mathtime}
% Note that miktex, by default, configures the mathtime package to use commercial fonts
% which you may not have. If you would like to use mathtime but you are seeing error
% messages about missing fonts (mtex.pfb, mtsy.pfb, or rmtmi.pfb) then please see
% the technical support document at http://www.aeaweb.org/templates/technical_support.pdf
% for instructions on fixing this problem.

% Note: you may use either harvard or natbib (but not both) to provide a wider
% variety of citation commands than latex supports natively. See below.

% Uncomment the next line to use the natbib package with bibtex
\usepackage{natbib}

% Uncomment the next line to use the harvard package with bibtex
%\usepackage[abbr]{harvard}

% This command determines the leading (vertical space between lines) in draft mode
% with 1.5 corresponding to "double" spacing.
\draftSpacing{1.5}


% tightlist command for lists without linebreak
\providecommand{\tightlist}{%
  \setlength{\itemsep}{0pt}\setlength{\parskip}{0pt}}



\usepackage{pdfpages}
\usepackage{graphicx}
\usepackage{caption}
\usepackage{booktabs}
\usepackage{threeparttable} % for table notes
\usepackage{caption}

\usepackage{hyperref}

\begin{document}


\title{The Effect of Airbnb on Long-Term Housing Availability and
Prices}
\shortTitle{The effect of Airbnb on the housing market}
% \author{Author1 and Author2\thanks{Surname1: affiliation1, address1, email1.
% Surname2: affiliation2, address2, email2. Acknowledgements}}


\author{
  Luca Gisler\thanks{
  : University of Science and Arts of Oklahoma, \href{mailto:}{}.
}
}

\date{\today}
\pubMonth{01}
\pubYear{2026}
\pubVolume{}
\pubIssue{}
\JEL{}
\Keywords{Airbnb, Housing Affordability}

\begin{abstract}
This paper empirically estimates how Airbnb activity is associated with
changes in housing affordability and availability. It provides a rough
estimate of the average increase in long-term rental apartment prices by
the number of Airbnb listings in an area. A 1\% increase in Airbnb
listings in an area in 2025 is associated with a X\% increase in
long-term rental prices at both 10\% and 5\% levels of significance.
\end{abstract}


\maketitle

According to Jarvis (2024), the Dallas--Fort Worth metropolitan area is
among the fastest-growing regions in the United States. This rapid
growth places upward pressure on housing demand and prices. Dallas--Fort
Worth has also experienced one of the highest growth rates of Airbnb
listings in the United States (CultureMap Dallas).The spatial mismatch
hypothesis highlights the close relationship between housing and labor
markets. Employment opportunities are often concentrated near city
centers, while affordable housing is increasingly located farther away.
This separation leads to longer commute times and reduced job
accessibility for low-income households (Ihlanfeldt).As tourism in the
Dallas--Fort Worth area has increased (Graber), short-term rental
platforms such as Airbnb have expanded rapidly. While these platforms
may increase income opportunities for property owners, they can also
incentivize the conversion of long-term rental units into short-term
accommodations. This study examines whether such conversions reduce
long-term housing availability and contribute to rising housing costs in
Dallas.

\subsection{Background and Motivation}\label{background-and-motivation}

Housing is a fundamental necessity, yet many households are unable to
afford homeownership and must rely on rental markets. This research adds
to prior work on homelessness, where housing affordability is a key
factor preventing homelessness. In urban areas with constrained housing
supply, competition for rental units can drive inequality by increasing
commuting distances and limiting access to employment opportunities.

\subsection{Research Question and
Limitations}\label{research-question-and-limitations}

This research examines whether increased participation in short-term
rental platforms contributes to declining housing affordability in
Dallas. Google Trends data and Airbnb listings are used as proxies for
short-term rental demand. Key limitations include potential measurement
error, seasonal variation in tourism, and proprietary Airbnb data.

\subsection{Economic Significance and
Application}\label{economic-significance-and-application}

Understanding how short-term rentals affect housing supply is crucial
for policymakers. If platforms such as Airbnb reduce the availability of
long-term housing, targeted regulations or zoning policies may be
required to preserve affordability and mitigate spatial inequality.

\section{Literature Review}\label{literature-review}

\subsection{Two-sided Market and Market
Efficiency}\label{two-sided-market-and-market-efficiency}

According to (Rysman, 2009) a market is Two-sided if there are two
agents interacting through an intermediary or a platform, while each
decision made by an Agent directly influences the outcome of the
opposite agent. While some One-sided markets also have an intermediary
or a platform the main distinction often is that the decision made by
each agent don't influence the outcome of the other agent. A good
example for a One-sided market with a platform is Netflix, where the
platform connects the viewers with Movies/Shows but the decisions made
by one viewer doesn't directly influence another agent.

Airbnb is a very good example for a Two-sided Market as the platform
connects hosts, who list their place for short-term rent, with Guests
that are looking to find a home for short-term rental. The agents in the
case of Airbnb are the Hosts and Guests. The distinction between a
one-sided market is that both agents directly influence the outcome for
eachother. For Example if there are more hosts offering places, it may
result in potential lower prices for prices for guests as there are more
options. On the other side more guests will result in higher occupancy
and therefore higher revenue for the hosts, this principle is often
referred to as cross-side network effects. According to (Holmgren, 2023)
platforms like AirBnb have to come up with strategies to find the
equilibrium of Hosts and Guests in order to maximize the platform
efficiency. Too Many listings might not meet enough demand and if there
are too many travelers there is a high chance that Airbnb listings are
already sold out. Therefore the way the two agents interact with each
other can have both positive and negative cross-side network effects.

\subsection{Market Efficency}\label{market-efficency}

As explained in the previous subsection platforms like Airbnb often
choose strategies that help the Ecosystem to restore the Equilibrium. In
order to understand Airbnb it is important to see the benefits each
agent has over traditional short-term rental models like Hotels. The
main advantage for hosts is that they can rent out houses that are
seasonally vacant or rooms in their house that are unused to guests
through a simple platform. Landlord also benefit from plattforms like
Airbnb as studies like (Hill et al., 2023) have found that the average
Landlord receives almost twice the long-term rental rate when using
short-term rental platforms. On the other side guests benefit from a
larger selection of price classes and locations compared to traditional
hotels. Guest also have the unique opportunity to find good deals
through dynamic pricing during off-seasons. These examples illustrate
the efficiency gains from short-term rentals on both the supply and
demand side.

\subsection{Short-term and long-term Rental
Markets}\label{short-term-and-long-term-rental-markets}

\begin{figure}[h] % h = "here" in text
    \centering
    \includegraphics[width=0.9\textwidth]{/Users/luca/Desktop/Housingvacancy_dallas_map.pdf}
    \caption{Housing Vacancy Map — Dallas, TX.}
    \label{fig:dallasmap1}
\end{figure}

\begin{figure}[h] % h = "here" in text
    \centering
    \includegraphics[width=0.9\textwidth]{/Users/luca/Desktop/Screenshot 2026-01-19 at 7.57.57 PM.png}
    \caption{Spatial Distribution of Airbnb Listings in Dallas County}
    \label{fig:dallasmap2}
\end{figure}

\subsection{Empirical Evidence on Airbnb's Effect on Housing Supply and
Rents}\label{empirical-evidence-on-airbnbs-effect-on-housing-supply-and-rents}

\subsection{Gaps in Existing
Literature}\label{gaps-in-existing-literature}

\section{Theoretical Framework}\label{theoretical-framework}

\subsection{Conceptualizing Airbnb as Two-Sided
Platform}\label{conceptualizing-airbnb-as-two-sided-platform}

\subsection{Market Equilibrium Between Short-term and Long-term
Rentals}\label{market-equilibrium-between-short-term-and-long-term-rentals}

As more hosts join the platform, housing units are increasingly
reallocated to the short-term market, leading to a growing number of
guests that can benefit from a greater variety of accommodation options.

\subsection{Expected Effects on Housing Supply and
Prices}\label{expected-effects-on-housing-supply-and-prices}

\section{Data and Methodology}\label{data-and-methodology}

\subsection{Dependant Variable}\label{dependant-variable}

This research paper's dependent variable is the demand for short-term
rentals in Dallas, mea- sured by Google Trends and the number of active
Airbnb listings in the city. Since the primary focus of this study is to
explore the additional factors contributing to the housing affordabil-
ity crisis, the rent-to-income ratio was used as a key measure of
housing affordability in the analysis. For more reliable statistics, the
data also included the U.S. Census Bureau's Ameri- can Community Survey
(ACS) one-year estimates, as well as Federal Reserve Economic Data
(FRED), to measure additional key economic indicators such as median
rent, household in- The Impact of Airbnb on Housing Supply and Prices 5
come, and population estimates. This approach enabled more
comprehensive, statistically rel- evant measurement of short-term rental
demand in Dallas and of housing affordability.

\#Dependant Variable is renter and renter units

\subsection{Independant Variable}\label{independant-variable}

The primary objective of this study is to examine the relationship
between the usage of plat- forms such as Airbnb and Housing
Affordability. Due to the difficulty of directly measuring short-term
rental demand, this analysis uses Google Trends data points and Airbnb
Listings as proxies for it.

\#Airbnb Density

\subsection{Variables and Summary
Statistics}\label{variables-and-summary-statistics}

\subsection{Econometric Model
specification}\label{econometric-model-specification}

\subsection{Limitations}\label{limitations}

\section{Discussion}\label{discussion}

\subsection{Regression Results and
Discussion}\label{regression-results-and-discussion}

The first table is looking at the effects of Airbnb density on the cost
of rent in short-rental apartments across from Dallas. The number of
observations in this case represent the amount of census tracts that
were used for the regression-line given that their population wasn't
zero, they didn't have zero renter units and their reported median rent
was valid. Other variables such as income, renter share and population
were held constant in accordance to the principle of ceteris paribus.
The initial results from that regression analysis suggest that there is
a negative correlation between Airbnb Density and long-term housing
rentals. The Analysis plot with the Regression Analysis line drawn
showed that the patterns changed in renter-heavy neighboorhoods.
Therefore the second model is looking at whether the Airbnb effects are
dependent on the amount of renter in the census tract.Because this paper
is looking at a potential housing supply shocked caused by Airbnb, the
third model switches from using the dependant variable of average rent
to available housing rental units. This model investigates whether
Airbnb density is associated with changes in the supply of
rental-appartments. In question 4 we are asking again if the housing
supply changes in areas, where there are many people renting
appartments. \begingroup \centering

\begin{tabular}{lcccc}
   \tabularnewline \midrule \midrule
   Dependent Variables: & \multicolumn{2}{c}{log\_median\_rent} & \multicolumn{2}{c}{ln\_renter\_units}\\
    & \multicolumn{2}{c}{Rent (log median rent)} & \multicolumn{2}{c}{Supply (log renter units)} \\ 
   Model:                     & (1)             & (2)             & (3)             & (4)\\  
   \midrule
   \emph{Variables}\\
   Constant                   & 2.520$^{***}$   & 2.755$^{***}$   & 12.52$^{***}$   & -1.435$^{**}$\\   
                              & (0.3867)        & (0.3791)        & (1.093)         & (0.5668)\\   
   Log Airbnb Density         & -0.0207$^{***}$ & -0.0454$^{***}$ & 0.2302$^{***}$  & 0.1155$^{***}$\\   
                              & (0.0075)        & (0.0091)        & (0.0309)        & (0.0167)\\   
   Log Median Income          & 0.4352$^{***}$  & 0.4268$^{***}$  & -0.8691$^{***}$ & 0.0185\\   
                              & (0.0245)        & (0.0243)        & (0.0713)        & (0.0382)\\   
   Renter Share (centered)    & 0.0160          & -0.1021$^{***}$ &                 & 2.885$^{***}$\\   
                              & (0.0309)        & (0.0338)        &                 & (0.0646)\\   
   Log Population             & -0.0118         & -0.0283         & 0.4094$^{***}$  & 0.9101$^{***}$\\   
                              & (0.0219)        & (0.0216)        & (0.0860)        & (0.0400)\\   
   airbnb / Renter            & 0.1837$^{***}$  &                 &   \\   
                              &                 & (0.0297)        &                 &   \\   
   Airbnb / Renter Share      &                 &                 &                 & -0.1344$^{**}$\\   
                              &                 &                 &                 & (0.0574)\\   
   \midrule
   \emph{Fit statistics}\\
   Observations               & 617             & 617             & 617             & 617\\  
   R$^2$                      & 0.53016         & 0.55853         & 0.29626         & 0.87406\\  
   Adjusted R$^2$             & 0.52709         & 0.55492         & 0.29281         & 0.87302\\  
   \midrule \midrule
   \multicolumn{5}{l}{\emph{Clustered (GEOID) standard-errors in parentheses}}\\
   \multicolumn{5}{l}{\emph{Signif. Codes: ***: 0.01, **: 0.05, *: 0.1}}\\
\end{tabular}
\par\endgroup

\subsection{Robustness Checks}\label{robustness-checks}

\subsection{Discussion of Findings}\label{discussion-of-findings}

\section{Conclusion and
Recommendations}\label{conclusion-and-recommendations}

\subsection{Summary of Key Findings}\label{summary-of-key-findings}

\subsection{Limitations and Future
Research}\label{limitations-and-future-research}


\end{document}
